\documentclass[11pt]{article}

\usepackage{amsmath,amssymb}
\usepackage{xurl}
\usepackage[hidelinks]{hyperref}
\setlength{\emergencystretch}{2em}

% Shared commands for notes in docs/paper-gaps/.
%
% Two halves.  The link commands below are project identity: OWNER/REPO is the
% placeholder that scripts/bootstrap_project.py replaces with this project's
% GitHub slug and Pages site.  Everything after them is NOTATION, and notation
% belongs to the source paper: the definitions here are an example set, and a
% project replaces them with the macros of its own paper's style file.  The one
% rule that never changes: a command defined here must mean exactly what the
% source paper's macro means, or a note silently says something else.

\newcommand{\Lean}{\textsc{Lean}}
\newcommand{\leanid}[1]{\textsf{#1}}
\newcommand{\ghissue}[1]{\href{https://github.com/OWNER/REPO/issues/#1}{GitHub issue \##1}}
\newcommand{\ghpr}[1]{\href{https://github.com/OWNER/REPO/pull/#1}{PR \##1}}
% Local issue tracker (issues/NNNN-*.md in this repository); no remote URL.
\newcommand{\localissue}[1]{issue \##1 (\path{issues/})}
\newcommand{\doclink}[1]{\href{https://OWNER.github.io/REPO/docs/#1.html}{\path{#1}}}

% --- Notation: replace with the source paper's own macros --------------------
% \F, \N, \C, \R, \tr, \ot, \eps, \E, \bx, \bu, \bv, \by

\newcommand{\F}{\mathbb{F}}
\newcommand{\N}{\mathbb{N}}
\newcommand{\C}{\mathbb{C}}
\newcommand{\R}{\mathbb{R}}
\newcommand{\tr}{\mathrm{tr}}
\newcommand{\eps}{\epsilon}
\newcommand{\ot}{\otimes}
\newcommand{\E}{\mathop{\bf E\/}}

% bold random variables
\newcommand{\bu}{\boldsymbol{u}}
\newcommand{\bv}{\boldsymbol{v}}
\newcommand{\bx}{{\boldsymbol{x}}}
\newcommand{\by}{\boldsymbol{y}}

% T1 so literal < and > in \texttt placeholders typeset as themselves
\usepackage[T1]{fontenc}

% bra-ket
\usepackage{braket}

% --- Example structured spaces ---
\newcommand{\polyfunc}[3]{\mathcal{P}(#1, #2, #3)}
\newcommand{\polymeas}[3]{\mathrm{PolyMeas}(#1, #2, #3)}
\newcommand{\polysub}[3]{\mathrm{PolySub}(#1, #2, #3)}

% Approximation relations (paper uses \approx_\eps and \simeq_\zeta)
\newcommand{\approxeps}{\approx_{\varepsilon}}
\newcommand{\simgeq}{\simeq_{\zeta}}


\title{Dimension Divisibility in the Classical Test Instantiation of the
	Pauli Basis Analysis}
\author{}
\date{2026-09-23}

\begin{document}
\maketitle

\section{At a glance}

\begin{itemize}
  \item \textbf{Difficulty: medium for the source comparison.} The printed
  analysis invokes a seed-indexed game at dimension $2m+2$, requiring
  $2m+2\mid q$. This fails for every admissible $m\ge2$ and for
  $(q,m)=(2,1)$; the exception is $m=1$, $q\ge8$. The formalization uses a
  directly indexed auxiliary game at that dimension. Its soundness, the
  seed-indexed soundness theorem for completed point POVMs on its original
  domain, and the final
  polynomial-pair conclusion are proved by the low individual degree route.
  The printed tensor-code reduction and the intermediate source-law
  comparisons are not thereby proved. The prescribed-effect statement is
  also distinguished from the completed-POVM statement linked in Lean.
  \item \textbf{Estimated weight: documentation only.} This revision
  compares existing definitions and results; it adds no proof or theorem
  hypothesis. The verdict is a \emph{documented deviation}: the auxiliary
  replacement is justified for the proved conclusions, while the precise
  limits of its correspondence with the source remain explicit. This is
  the documentary closure authorized for \ghissue{710}, not a claim that
  the unresolved intermediate assertions are proved or false.
  \item \textbf{Mathlib/project split:} the existing argument is principally
  project-local. Mathlib supplies finite-field, matrix, and finite-sum
  infrastructure; the game transports, combining reduction, and polynomial
  agreement estimates are established in the repository.
  \item \textbf{Key mathematical inputs:} balanced finite sampling, the low
  individual degree theorem, correlated seed registers, and
  Schwartz--Zippel. The two separate tensor-code import obligations are
  retained in Section~\ref{sec:import-obligations} (\ghissue{527}); the
  original investigation is \localissue{0002}.
\end{itemize}

\section{Key theorem forms}

Throughout, $q = 2^r$ with $r$ positive and odd is an admissible field size.
The source's definition of an admissible tuple adds $m\mid q$; the
formal parameter records make positivity of the dimension, degree, and
simultaneity explicit. We work with these positive parameters, and
$\mathrm{ideg}_{d,m}(\F_q)$ denotes the $m$-variate polynomial representatives
over $\F_q$ of individual degree at most $d$.

The classical soundness theorem asserts that, for $m'\mid q$, $\eps>0$,
and a projective strategy of value at least $1-\eps$ for the seed-indexed
game at $(q,m',d,k)$, there are polynomial-tuple POVMs $G^{\mathrm A}$ and
$G^{\mathrm B}$ on the original players' spaces such that
\[
 A^u\simeq_{\delta_{\rm ld}}G^{\mathrm B}(u),\qquad
 G^{\mathrm A}(u)\simeq_{\delta_{\rm ld}}B^u,\qquad
 G^{\mathrm A}\simeq_{\delta_{\rm ld}}G^{\mathrm B},
 \tag{\path{lem:ld-soundness}}
\]
where the first two relations average over uniform $u\in\F_q^{m'}$, all
relations use the strategy state and opposite tensor placements, and
$\delta_{\rm ld}=a(dm'k)^a(\eps^b+q^{-b}+2^{-bm'd})$ for universal
$a\ge1$, $0<b\le1$. Here $G(u)$ means evaluation of every polynomial in
the outcome tuple. In this source display, $A^u_{\mathbf a}$ and
$B^u_{\mathbf a}$ are the strategy's effects for the prescribed point
answer $\mathbf a\in\F_q^k$.

The Lean theorem instead compares completed point POVMs. The strategy's
answer alphabet also contains axis-line and diagonal-line polynomial
answers. At a point question their effects have positive sums
\[
 R_A^u=I_{\mathrm A}-\sum_{\mathbf a\in\F_q^k}A^u_{\mathbf a},
 \qquad
 R_B^u=I_{\mathrm B}-\sum_{\mathbf a\in\F_q^k}B^u_{\mathbf a}.
\]
The point readout sends all these wrong-form answers to the zero tuple, so
the complete families are
\[
 \widetilde A^u_{\mathbf a}
   =A^u_{\mathbf a}+\mathbf1_{\mathbf a=\mathbf0}R_A^u,
 \qquad
 \widetilde B^u_{\mathbf a}
   =B^u_{\mathbf a}+\mathbf1_{\mathbf a=\mathbf0}R_B^u.
\]
The established statement has the same hypotheses, error function,
question laws, and original players' spaces, but its first two relations are
\[
 \widetilde A^u\simeq_{\delta_{\rm ld}}G^{\mathrm B}(u),
 \qquad
 G^{\mathrm A}(u)\simeq_{\delta_{\rm ld}}\widetilde B^u.
\]
It also proves the unchanged global relation
$G^{\mathrm A}\simeq_{\delta_{\rm ld}}G^{\mathrm B}$, for every positive
$k$, with no tensor-code correspondence or parameter-bound
hypothesis.\footnote{\leanid{exists\_ld\_soundness},
\doclink{MIPStarRE/QPBT/Test/LowDegreeGameTheorems.lean}, lines~82--112;
\leanid{LdParams}, \doclink{MIPStarRE/QPBT/Test/LowDegreeGame.lean},
lines~40--50; \leanid{ldPointValuesOrZero},
\doclink{MIPStarRE/QPBT/Test/LowDegreeGameMeasurements.lean}, line~370.}
No hypothesis forces $R_A^u$ or $R_B^u$ to vanish. The two theorem displays
therefore do not agree verbatim. Their mathematical relation follows from
positivity: the consistency defect is an average of
$\sum_{\mathbf a\ne\mathbf b}
\langle\psi|A^u_{\mathbf a}\ot C^u_{\mathbf b}|\psi\rangle$,
where $C^u$ is the other player's positive measurement family.
Replacing $A^u$ by $\widetilde A^u$ adds only nonnegative summands,
and the same applies to $B^u$. Hence the completed bounds imply the
prescribed-effect bounds with the same error. This paragraph explains the
comparison in ordinary mathematics; the linked Lean signature does not
state that implication, and no new kernel proof is supplied here. The
source-labelled blueprint entry remains unmarked with a reasoned C4
exemption; the separately labelled completed statement retains its marks.

The original domain $m'\mid q$ must also be distinguished from the
extended-dimension application, whose comparison has four forms.
\begin{enumerate}
  \item \textbf{Source theorem form} (the polynomial-pair measurement).  There
  is a function $\delta_S(\eps,m,d,q) = a(md)^a(\eps^b + q^{-b} + 2^{-bmd})$,
  for universal constants $a > 1$ and $0 < b < 1$, such that for every
  admissible $(q,m,d)$ and every projective strategy for the Pauli basis test
  succeeding with probability at least $1-\eps$, there is a projective
  measurement $\{\hat S_{g_X,g_Z}\}$ with outcomes pairs of polynomials in
  $\mathrm{ideg}_{d,m}(\F_q)$ whose evaluations are $\delta_S$-consistent, on
  average over a uniform point $u \in \F_q^m$ and in both bases
  $W \in \{X,Z\}$, with the corresponding points measurements
  $\{\hat M^{(\mathrm{Point},W),u}_a\}$ of the strategy, together with the
  symmetric counterpart of this relation.
  \item \textbf{Missing instantiation obligation.}  The classical low
  individual degree game with parameter tuple $(q,\, 2m+2,\, d,\, 1)$ is well
  defined, and the combined point and line measurements assembled in the
  analysis constitute a strategy for it of value at least
  $1 - \delta_{\mathrm{combine}}$.  The game's question distribution draws
  its axis and diagonal indices through an index function
  $\chi_{2m+2}\colon \F_q \to \{1,\ldots,2m+2\}$ with classes of equal size
  $q/(2m+2)$, so the game exists only when $2m+2$ divides $q$.
  \item \textbf{Divisibility-restricted source candidate.}  The source theorem
  form under the additional hypothesis $2m+2 \mid q$.  Among
  admissible tuples this hypothesis holds exactly when $m = 1$ and
  $q \ge 8$; it fails for every admissible tuple with $m \ge 2$ and for
  $(q,m) = (2,1)$.  This restriction makes the auxiliary source game defined;
  it does not settle either soundness-import obligation, so this form is not
  established by the present route.
  \item \textbf{Established directly indexed form.}  Quantum
  soundness of the directly indexed variant of the classical low individual
  degree game over $\F_q^{m'}$, for every dimension $m' \ge 1$ and
  admissible $q$, with no divisibility hypothesis: the question
  distribution samples the axis index, and the diagonal prefix index,
  uniformly from $\{1,\ldots,m'\}$ directly rather than through an index
  function on $\F_q$, and line questions carry the sampled index itself in
  place of the seed $s \in \F_q$; for simultaneity $\kappa$, the error is
  $\delta_{\mathrm{ld}}(\eps,q,m',d,\kappa)
  = a(dm'\kappa)^a(\eps^b + q^{-b} + 2^{-bm'd})$.  This form is proved by
  applying the low individual degree theorem to a combined strategy at
  dimension $M=m'+\kappa$, with sampling count $N=2560000M^3d$, and then
  recovering the $\kappa$ polynomial coordinates.  The proof verifies the low
  individual degree theorem's condition $400Md \le N$; it does not invoke
  Theorem~4.7 of \cite{Ji2021Quantum} or verify that theorem's distinct
  condition $K \ge 12m'(d+1)$.  Thus it avoids rather than discharges the
  two source-import obligations of Section~\ref{sec:import-obligations}.
  Its point families likewise use the completed readout above.
  When $m'\mid q$, for every $k$, transport to the seed-indexed game is
  established: the two mixed relations compress exactly and the global relation is recovered
  from point agreement on the original Hilbert spaces and the
  Schwartz--Zippel bound for polynomial tuples.
\end{enumerate}

\section{The source assertion and the obstruction}

We examine the final step of the subsection ``Applying the classical
low-degree test'' in the analysis of the Pauli basis test in the compression
paper of Ji, Natarajan, Vidick, Wright, and Yuen (MIP$^{*}$=RE,
arXiv:2001.04383) \cite{Ji2020MIPStar}, whose quantum soundness component builds on the low
individual degree test of \cite{Ji2020LowIndividualDegree}.%
\footnote{Tracked in \localissue{0002}.  In the mirror of the source
consulted here, the polynomial-pair lemma and its proof are at
\path{references/qpbt-paper/14_analysis_of_the_pauli_basis_test.tex},
lines~1267--1288; the sub-line lemma is at lines~1063--1116 and the combined
line lemma at lines~1020--1034 of the same file.  The game parametrization
with its divisibility, the index function, the soundness statement with its
proof, and admissibility are at
\path{references/qpbt-paper/08_classical_and_quantum_low_degree_tests.tex},
lines~33--35, 214--219, 413--458, and 958--961.}
There, the combined point measurements
$\{\hat Q^{x,z,\alpha,\beta}_c\}$ and combined line measurements
$\{\hat Q^{\ell}_f\}$, indexed by points of $\F_q^{2m+2}$ and lines in
$\F_q^{2m+2}$, are asserted to form a strategy for the classical low
individual degree game with parameters
\[
  \mathsf{ldparams} = (q,\, 2m+2,\, d,\, 1)
  \tag{\path{lem:qld-4-7}}
\]
of value $1 - \delta_{\mathrm{combine}}$, and the quantum soundness theorem
for that game is applied to produce the global polynomial-pair measurement.
The two lemmas feeding this step are stated for the line-point distribution
over $\F_q^{2m+2}$.

The game, however, is parametrized by tuples $(q,m',d,k)$ with $q$ an
admissible field size \emph{and $m'$ dividing $q$}.  The divisibility is not
incidental: the question distribution is a conditionally linear distribution
in which a uniform seed coordinate $s \in \F_q$, read as an integer in
$\{0,\ldots,q-1\}$, determines the axis or diagonal index through the unique
$\chi(s)$ with
\[
  s = (\chi(s)-1)\,\frac{q}{m'} + r, \qquad 0 \le r < \frac{q}{m'},
  \tag{\path{eq:chi-func}}
\]
so that the $m'$ classes $\chi^{-1}(i)$ partition $\F_q$ into sets of equal
size $q/m'$ and the index is uniform. Indeed, each index has probability
$|\chi^{-1}(i)|/q$, which equals $1/m'$ only if every class has the integer
size $q/m'$. Thus for $m'\nmid q$ this balanced construction is unavailable,
and the source has not defined its game or its line-point law there. An
arbitrary map with unequal fibres would specify a different sampling law.

Admissibility of $(q,m,d)$ supplies only $m \mid q$.  Since an admissible $q$
is a power of two, $m \mid q$ forces $m = 2^j$.  For $j \ge 1$ the integer
$m+1 \ge 3$ is odd, so $2m+2 = 2(m+1)$ has an odd factor exceeding $1$ and
divides no power of two; for $m = 1$ the condition $4 \mid q$ holds for every
admissible $q$ except $q = 2$.  Thus the instantiation at dimension $2m+2$ is
undefined for every admissible tuple with $m \ge 2$: for instance
$(q,m,d) = (8,2,1)$ is admissible, while $6 \nmid 8$.  The source performs
the instantiation without addressing the divisibility.  Thus even before the
two soundness-import obligations below are considered, this route can apply to
the polynomial-pair lemma only for $m = 1$ and $q \ge 8$.

\section{What the cited soundness theorem requires: the import obligations}
\label{sec:import-obligations}

The quantum soundness theorem for the classical game inherits the constraint
$m' \mid q$ solely because it is stated for the game just described; its
proof fixes an admissible $q$ with $m' \mid q$, asserts that the game with
simultaneity parameter $1$ is \emph{identical} to the two-prover tensor-code
test for the Reed--Solomon tensor code $\mathcal C^{\ot m'}$ of the same
authors' \emph{Quantum soundness of testing tensor codes} (arXiv:2111.08131)
\cite{Ji2021Quantum}, and concludes by Theorem~4.7 there, instantiated with
that theorem's auxiliary integer parameter --- denoted $k$ in that source,
and written $K$ here to avoid clashing with the field exponent and the
simultaneity parameter --- chosen as $(m')^3 d$.  Neither
the asserted identity nor the instantiation is complete as stated, so the
import carries two verification obligations, which this section details and
which the blueprint's proof node for the soundness statement records as
open for the printed tensor-code route. Both obligations are independent of
the divisibility question and remain open for that route at every dimension
where that game is defined.  They would also arise from an attempt to prove
soundness of the directly indexed game by the same tensor-code import.  The
established proof for the directly indexed game takes a different route through
the low individual degree theorem and therefore bypasses both obligations.

\paragraph{The cited test.}
In the two-prover tensor-code test (Figures~1 and~2 of
\cite{Ji2021Quantum}; the two-prover strategy format and the goodness
measure are Definitions~4.5 and~4.6 there), the referee performs one of
three subtests with probability $1/3$ each, assigning the two roles of each
subtest to the two provers uniformly at random.  \emph{Axis-parallel lines
test}: sample a uniform point $u \in [n]^{m'}$ and a uniform axis
$j \in \{1,\ldots,m'\}$; send the axis-parallel line $\ell$ through $u$ in
direction $j$ to one prover, who answers a codeword $g \in \mathcal C$;
send $u$ to the other, who answers a field value $a$; accept iff
$g(u_j) = a$.  \emph{Subcube commutation test}: sample
$j \in \{1,\ldots,m'\}$ and a subcube $H \subseteq [n]^{m'}$ by fixing the
last $j-1$ coordinates uniformly; sample $u, v$ independently and uniformly
in $H$ and a uniform order bit; send the ordered pair $(u,v)$ or $(v,u)$ to
one prover, who answers a pair $(b_u, b_v)$ of field values; send $u$ to
the other, who answers $a$; accept iff $b_u = a$.  \emph{Synchronicity
test}: send one identical question --- a line, a point, or an ordered
pair, with probability $1/3$ each --- to both provers and accept iff their
answers agree.  Theorem~4.7 supposes a two-prover strategy passing this
test with probability $1-\eps$ together with an integer parameter
$K \ge 12 m' t$, where $t$ is the interpolation parameter of the base code
--- for the degree-$d$ Reed--Solomon code with $d<q$, of blocklength $n = q$ and
distance $q-d$, one has $t = d+1$ --- and produces projective measurements
$\{G_c\}$ and $\{\tilde G_c\}$ with outcomes in $\mathcal C^{\ot m'}$,
consistent with one prover's points measurements, with the other prover's
lines measurements, and with each other, up to
$\eta'(m',t,K,\eps,n^{-1}) = \mathrm{poly}(m',t,K) \cdot
\mathrm{poly}(\eps, n^{-1}, e^{-K/(m')^2})$.

\paragraph{Obligation 1: the game correspondence.}
The source asserts that the $\chi$-based game at parameters $(q,m',d,1)$
is this test.  It is not, in four respects.
(i)~The game has no pair questions.  A pair question of the test reveals
two points of a \emph{hidden} subcube --- the subcube itself is not sent
--- and is answered by two field values, while the game's diagonal-line
question reveals a full line description and is answered by a degree-$m'd$
polynomial.  A reduction must construct the pair measurements from the
diagonal-line measurements by post-processing (measure a line through $u$
and $v$ and evaluate the answered polynomial at the two parameters), and
both directions of the information mismatch matter: the game's line
description carries an index and a seed, neither of which is a function of
the pair $(u,v)$, so the constructed measurement must average over the
compatible indices and seeds; and the game's diagonal lines fix the
\emph{first} $i-1$ coordinates where the test's subcubes fix the
\emph{last} $j-1$, so the coordinate conventions must be aligned.  The
construction must justify the measurement class used in the import and must
match the law of the hidden-subcube pair with the law induced by the game's
diagonal sampling.
(ii)~The game's line questions carry the seed $s \in \F_q$; the test's carry
only the line.  A strategy may depend on the seed beyond its class
$\chi^{-1}(i)$.  For the comparison with the directly indexed game, this
dependence is now retained in a correlated residue register rather than
averaged away, as explained below.  That construction resolves this part of
the direct-game transport, but it has not been composed with the distinct
pair-question construction required for the tensor-code test.
(iii)~The branch weights differ: each nontrivial cross-test has mass $1/3$
in the test against $2/9$ in the game.  Once the per-branch predicates are
matched, the intended bookkeeping bound is that the constructed strategy's
rejection probability in the test is at most $\tfrac32$ times the game's
rejection probability, up to the $O(1/(m'q))$ of item~(iv); this bound is
the target of the obligation, not an established fact.
(iv)~The game's predicate fixes an acceptance convention on diagonal
questions with zero direction --- an event of probability at most
$2/(m'q)$, recorded in the blueprint's remark on zero directions --- for
which the test has no counterpart branch.
Finally, the displayed conclusions of Theorem~4.7 pair one prover's points
measurements with the other prover's tensor-codeword measurement, and the
first prover's tensor-codeword measurement with the second prover's lines
measurements, while the game-level statement asserts the point consistency
in both directions; deriving the stated form uses the role symmetrization
of the test and belongs to the correspondence.  None of these steps is
difficult to state, but none is carried out in the source, and together
they replace the asserted one-line identity by a reduction carrying its
own value bound.

\paragraph{Obligation 2: the parameter bound.}
Theorem~4.7 requires $K \ge 12 m' t = 12 m' (d+1)$.  The source
instantiates $K = (m')^3 d$, which meets the bound iff
$(m')^2 d \ge 12(d+1)$: this holds for all $d \ge 1$ when $m' \ge 5$,
holds at $m' = 4$ exactly when $d \ge 3$, and fails for every $d$ when
$m' \le 3$.  In the game's parameter domain the dimension divides the
admissible field size and is therefore a power of two, so the source's
instantiation violates the hypothesis exactly at $m' \in \{1,2\}$ for
every $d$ and at $m' = 4$ for $d \in \{1,2\}$ --- in particular, in the
multilinear regime $d = 1$, at every dimension $m' \le 4$.  A candidate
repair is to instantiate $K = \max((m')^3 d,\, 12 m'(d+1))$; the open
obligation is then to re-absorb the error, that is, to check that
$\mathrm{poly}(m',t,K)$ and $e^{-K/(m')^2}$ still collapse to the form
$a(dm')^a(\eps^b + q^{-b} + 2^{-bm'd})$ in the small-$m'$ cases where the
maximum is $12m'(d+1)$ --- plausible, since there
$e^{-K/(m')^2} = e^{-12(d+1)/m'}$ decays in $d$ at a rate comparable to
$2^{-m'd}$ for $m' \le 4$, but not checked here.

In the correspondence of Obligation~1, the only role of the divisibility
$m' \mid q$ is to make the $\chi$-sampled index uniform.  The underlying
mathematics does not need the divisibility.  The low
individual degree test of \cite{Ji2020LowIndividualDegree} samples its axis
index $i$ uniformly from $\{1,\ldots,m'\}$ directly, with no index function
on $\F_q$ and no divisibility hypothesis,%
\footnote{See the test figure in \path{references/ldt-paper/test_definition.tex}.}
and the tensor-code test is likewise a test about the code
$\mathcal C^{\ot m'}$, not about an encoding of indices into $\F_q$.  The
$\chi$-based packaging exists so that the question distribution is
conditionally linear, which the introspection machinery requires for the
game the players actually play --- the Pauli basis test at dimension $m$,
where admissibility does supply the divisibility.  The game at dimension
$2m+2$ never appears in a verifier; it is an internal analysis device in the
proof of the polynomial-pair lemma, so nothing there requires its question
distribution to be conditionally linear.

\section{Established formalization and alternatives}

\paragraph{Direct index sampling (soundness established).}
For arbitrary $m' \ge 1$, the formalization defines the line-point
distribution over $\F_q^{m'}$ by sampling the point, the axis index
$i \in \{1,\ldots,m'\}$, and the diagonal direction with vanishing first
$i-1$ coordinates directly.  In the resulting game, line questions carry
the index $i$ (and, for diagonal lines, the projected direction) in place
of the seed $s \in \F_q$.  Soundness form~4 is established for this game
through the low individual degree theorem and the general-simultaneity
combining reduction analysed in \cite{gap:qpbt_ld-simultaneous-sandwich}.
This removes the divisibility obstruction for the
auxiliary directly indexed game. The constants supplied by the general
simultaneous proof are $a=10^{23}$ and $b=1/80000$.\footnote{
\leanid{exists\_direct\_ld\_soundness},
\doclink{MIPStarRE/QPBT/Combining/DirectLowDegree/Soundness.lean}, line~65;
\leanid{exists\_directCombinedTransportConstants},
\doclink{MIPStarRE/QPBT/Combining/DirectLowDegree/Transport/Combining/Error.lean}.
The carrier and sampling definitions are \leanid{DirectLdParams},
\leanid{DirectLineDesc}, and \leanid{directLinePointDist} in
\doclink{MIPStarRE/QPBT/Combining/DirectLowDegree/Geometry.lean},
lines~36, 136, and 263.}
Because the proof does not use the
tensor-code import, it leaves both original import obligations open rather
than establishing them.

The relation to the present game when $m' \mid q$ must be proved at the level
of strategies, not merely of question distributions.  An
$\mathsf{ALine}$ question of the $\chi$-based game is a pair $(u_0,s)$ and a
$\mathsf{DLine}$ question is a triple $(u_0,s,v')$, so measurements may vary
between seeds in the same class $\chi^{-1}(i)$.  Each seed nevertheless has a
unique decomposition into its index $i$ and a residue
$\rho\in\{0,\ldots,q/m'-1\}$.  The seed-indexed question law is thereby
identified with the direct question law together with a uniform residue; the
parsing map back to the direct question law is an exact pushforward.

The formal transport retains, rather than averages away, the dependence on
$\rho$.  It adjoins a register of dimension $q/m'$ to each player's Hilbert
space, places the two registers in a maximally correlated state, and makes
each question measurement block diagonal, with the block at $\rho$ equal to
the original measurement for the corresponding seed.  Thus both players see
the same residue on a synchronicity branch.  This correlated-residue lift
preserves normalization and projectivity and preserves the game value exactly,
branch by branch.  In particular, seed-dependent relabellings cause no loss;
the earlier obstruction arose only from independently averaging the local
measurements, which is not the construction used here.%
\footnote{The exact distribution comparison is
\leanid{ldQuestionDistribution\_map\_parse} in
\path{MIPStarRE/QPBT/Combining/DirectLowDegree/Transport/Correspondence.lean},
line~389; it requires an \leanid{LdParams} record, including $m'\mid q$.
The lift, projectivity theorem, and value identity are
\leanid{ldStrategyToDirect},
\leanid{ldStrategyToDirect\_isProjective}, and
\leanid{ldStrategyToDirect\_value\_eq} in
\path{MIPStarRE/QPBT/Combining/DirectLowDegree/Transport/SeedFiber.lean} and
\path{MIPStarRE/QPBT/Combining/DirectLowDegree/Transport/SeedFiberValue.lean}.}

Applying form~4 to the lifted strategy produces polynomial-tuple POVMs
$G^{\mathrm A}$ and $G^{\mathrm B}$ on the enlarged carriers.  Compress one
such POVM by taking the normalized average of its diagonal residue blocks.
When the other family is a completed point measurement, it is constant on the opposite
residue register, and this compression preserves the associated defect
exactly.  Consequently the point/polynomial and polynomial/point conclusions
of the directly indexed theorem transport back to the original seed-indexed
strategy with no additional error.%
\footnote{These two identities are
\leanid{ldStrategyToDirect\_pointPolynomial\_compression} and
\leanid{ldStrategyToDirect\_polynomialPoint\_compression} in
\path{MIPStarRE/QPBT/Combining/DirectLowDegree/Transport/Consistency/Compression.lean}.}

The third conclusion is different.  On the lifted state, the
global/global defect pairs equal residues.  Compressing both polynomial
measurements separately instead produces a double average over independent
residues.  Neither quantity bounds the other in general.  For example, on a
two-element residue register, if both deterministic measurements report the
residue, the equal-residue disagreement is zero while the independent-residue
disagreement is $1/2$.  If Alice reports the residue and Bob reports its
complement, the corresponding disagreements are $1$ and $1/2$. For every $k$,
the third relation is instead recovered on the original Hilbert spaces.
The point-agreement branch has weight $1/9$, so the two point measurements
have consistency defect at most $9\eps$. The triangle inequality combines
this with the two compressed mixed relations to give agreement of evaluated
polynomial tuples. Distinct tuples differ in at least one coordinate, so
their evaluations can agree only where those two coordinate polynomials
agree. Schwartz--Zippel therefore gives the same collision bound $m'd/q$
for every $k$, without a union bound over the coordinates.
Writing $E=\delta_{\mathrm{ld}}(\eps,q,m',d,k)$ for the direct error, the
global defect after compression is at most
\[
  E+2\sqrt{9\eps+E}+\frac{m'd}{q}.
\]
For $\eps\leq1$ the direct error dominates both $\eps$ and $m'd/q$.
If also $E\leq1$, the displayed expression is at most $10\sqrt E$.
Replacing the direct constants $(a,b)$ by $(10a,b/2)$ absorbs this bound:
the square root of a sum is at most the sum of its square roots, and the
prefactor $a(dm'k)^a$ is at least one. When $E>1$ or $\eps>1$, the unit
bound on bipartite consistency defects suffices. This proves seed-indexed
soundness for completed point POVMs for every $k$ on its original domain
$m'\mid q$, with possible
universal constants $a=10^{24}$ and $b=1/160000$.%
\footnote{The point-agreement estimates are proved in
\doclink{MIPStarRE/QPBT/Combining/DirectLowDegree/Transport/PointAgreement.lean}.
The resulting theorem is \leanid{exists\_ld\_soundness} in
\doclink{MIPStarRE/QPBT/Test/LowDegreeGameTheorems.lean}; the scalar estimate is
\leanid{ten\_sqrt\_deltaLd\_le} in
\doclink{MIPStarRE/QPBT/Combining/DirectLowDegree/Transport/SeedError.lean}.}

The combining reduction has now been formalized with the direct distribution
and the directly indexed game, so it is well defined at dimension
$m'+\kappa$ for
all positive parameters.  Its proof uses the sampled geometric data and then
applies the low individual degree theorem at the combined dimension.  This
establishes soundness with completed point POVMs for the directly indexed
game and, by the compression and recovery argument above, for the
seed-indexed game. The directly indexed
application at dimension $2m+2$ now also constructs the global polynomial-pair
measurement, as \leanid{exists\_globalPairWitness} in
\path{MIPStarRE/QPBT/Combining/Apply.lean}. This proves the polynomial-pair
conclusion without a divisibility hypothesis. It does not identify the
undefined source seed law with the direct law or prove the intervening
source-shaped subline statements.
The external check that the tensor-code test of arXiv:2111.08131 samples
its axes and subcubes directly has been made against that source: Figure~1
there draws the axis index and the subcube data uniformly and directly,
with no index encoding into $\F_q$.

\paragraph{Padding the dimension.}
The following two alternatives record the earlier mathematical comparison;
neither is used by the established proof or adopted as a new game definition.
One could embed the combined polynomial into
$\mathrm{ideg}_{d,\bar m}(\F_q)$ for $\bar m = 4m$, the least power of two
with $\bar m \ge 2m+2$, ignoring the $2m-2$ dummy variables.  Then
$\bar m \mid q$ holds whenever $q \ge 4m$, and the existing game applies at
dimension $\bar m$.  In the complementary regime $m \in \{q/2,\, q\}$ the
conclusion is vacuous: $md \ge q/2$ gives
$\delta_S \ge a(md)^a q^{-b} \ge a\,2^{-a} q^{a-b} \ge a\,2^{-b} \ge 1$
for $a \ge 2$, while the consistency quantities bounded by the theorem never
exceed $1$.  This route stays inside the present definitions but reworks the
sub-line distribution (a fourth coordinate block of dummies), the combining
map, and the wrong-block Schwartz--Zippel analysis, and it moves the
combining lemmas away from their source statements at dimension $2m+2$.

\paragraph{Unequal index classes.}
One could instead define an index function at dimension $2m+2$ with classes
of sizes $\lfloor q/(2m+2)\rfloor$ and $\lceil q/(2m+2)\rceil$.  The index
is then only approximately uniform, the identification with the external
test fails exactly, and every distributional identity in the sub-line
analysis acquires an error term; this amounts to re-proving the imported
soundness theorem for a perturbed game and is not recommended.

\section{Scalar estimates on the auxiliary subline law}
\label{sec:auxiliary-subline-scalars}

The source sub-line lemma \path{lem:qld-sublines} samples the extended line
over $\F_q^{2m+2}$, so its distribution $D$ uses the seed-indexed
construction whose divisibility condition fails outside the exceptional
range $m=1$, $q\ge8$. The scalar
claims \path{claim:17-1} and \path{claim:17-3}, at lines 1140--1166 and
1204--1239 of
\path{references/qpbt-paper/14_analysis_of_the_pauli_basis_test.tex} --- the
blueprint nodes \path{lem:claim-17-1} and \path{lem:claim-17-3} --- are
stated over that $D$.  Their formal counterparts are stated over a different
law $D_{\rm dir}$, carried by \leanid{SubLineWitness} at line 209 of
\doclink{MIPStarRE/QPBT/Combining/Witnesses.lean}: its extended-line
marginal is the marginal of \leanid{directLinePointDist} (line 263 of
\doclink{MIPStarRE/QPBT/Combining/DirectLowDegree/Geometry.lean}), and its
$X$ and $Z$ point marginals are recorded \emph{separately} as mixtures of
the restricted line-point law \leanid{restrictedLinePointDist} (line 151 of
the same witnesses file), built from \leanid{restrictedALineDist} and
\leanid{restrictedDLineDist} at lines 379 and 387 of
\doclink{MIPStarRE/QPBT/Combining/Defs.lean}, which are the refined
restricted laws of blueprint \path{def:ith-restricted-line}.  That structure
claims no equality for the joint conditional law of the pair of projected
points.  Identifying the corresponding source push-forwards, including
forgetting the residual seed on axis lines, remains open.  Existence of this
auxiliary distribution does not identify it with $D$: sampling a coordinate
index directly does not construct the seed-indexed extended-line questions
at which the divisibility condition fails, which is the obstruction recorded
here.

\paragraph{The scalar conclusions.}
Let $A_Q(\cdot)$, $A_{ZX}(\cdot)$ and $A_Z(\cdot)$ denote the averaged
joint-point, ordered-product, and $Z$-point overlaps of these claims.  The
source assertions are
\[
 |A_Q(D)-A_{ZX}(D)|=O(\sqrt{\delta_Q}),\qquad
 |A_Z(D)-1|=O\bigl(\sqrt m\,
   (\delta_P^{1/4}+\delta_Q^{1/4}+\eps^{1/4})\bigr).
\]
The established statements instead bound real parts over $D_{\rm dir}$: for
projective joint point measurements and complete paired-line measurements
satisfying the witness hypotheses, and any auxiliary subline witness,
\begin{align*}
 |\operatorname{Re}A_Q(D_{\rm dir})-
     \operatorname{Re}A_{ZX}(D_{\rm dir})|&\le C\,\delta_Q^{1/2},\\
 |\operatorname{Re}A_Z(D_{\rm dir})-1|&\le
 C\sqrt m\,(\delta_P^{1/4}+\delta_Q^{1/4}+\eps^{1/4}).
\end{align*}
The constant is existentially quantified in both formal statements; the
recorded proofs supply the witness $C=2$, at lines 117 and 686 of
\doclink{MIPStarRE/QPBT/Combining/Claims.lean}.  The evaluations use the
completed answer alphabet, so an undefined evaluation contributes zero point
effect.  The joint-point and paired-line witnesses retain their stated
consistency and degree hypotheses; neither new transport data nor a reality
assumption is imposed on these results.

\paragraph{Which claims are real-part only.}
The ordered product $M_ZM_X$ need not be Hermitian, so a real-part bound
alone does not establish the first source assertion; the source
correspondence of the three scalar claims is kept in
\path{docs/paper-gaps/qpbt_subline-claims-line-marginal.tex}
(\ghissue{474}) and is not restated here.  The \path{lem:claim-17-2}
analogue over $D_{\rm dir}$ is not real-part only:
\leanid{subline\_remove\_X\_factor\_re\_direct} and
\leanid{subline\_remove\_X\_factor\_direct} in
\doclink{MIPStarRE/QPBT/Combining/Claims.lean} establish the real-part and
the complex-modulus form, the latter without a reality hypothesis.  The two
real-part estimates displayed above are
\leanid{subline\_replace\_by\_ordered\_product\_re\_direct} and
\leanid{subline\_Z\_term\_near\_one\_re\_direct} in the same file.

\paragraph{Remaining obligations.}
Two obligations with distinct mathematical content remain for those two
claims; neither is declared in the tree in any form.  A source-distribution
transport must construct the seed-indexed law where it is defined and prove
the corresponding push-forward, evaluation, and separate point-marginal
identities; over the unrestricted source domain it still requires the
definition-level obstruction above to be resolved.  A complex-modulus
estimate for the ordered-product claim must be derived from positivity,
completeness, the point marginal, and the same-placement distance, without
an added reality hypothesis.  For the $Z$ overlap, positivity of the two
commuting opposite-placement factors gives reality, which must also be
recorded when passing from the real-part result to the source scalar
statement.  These are described future proof targets, not assumed or
axiomatized declarations.

\paragraph{Blueprint bookkeeping.}
The blueprint therefore leaves \path{lem:claim-17-1} and
\path{lem:claim-17-3} pending and records the established auxiliary results
as the separate nodes \path{lem:claim-17-1-re-direct},
\path{lem:claim-17-2-direct-real}, \path{lem:claim-17-2-direct} and
\path{lem:claim-17-3-re-direct} of
\path{blueprint/src/chapter/ch15_qpbt_combining.tex}, with their carrier,
completed evaluations, measurement hypotheses, and their real-part or
complex-modulus conclusions explicit; their Lean statements and proofs are retained.  This scope record
introduces no new proof hole, and it changes nothing about the source-import
obligations of Section~\ref{sec:import-obligations} (\ghissue{527}), nor
about the separate obligations tracked in
\path{docs/paper-gaps/qpbt_combined-lines-error-term.tex} (\ghissue{598})
and \path{docs/paper-gaps/qpbt_symmetrization-attainment.tex}
(\ghissue{524}).

\section{Comparison with the blueprint}

The blueprint \cite{MIPStarREBlueprint} keeps the sub-line lemma, the
combined line lemma, and the polynomial-pair lemma in their source-shaped
statements, with no added divisibility hypothesis. It records the failed
source instantiation at $(q,2m+2,d,1)$ and leaves the source sub-line and
combined-line assertions unmarked. In contrast, the polynomial-pair lemma
has a complete proof through the directly indexed construction, which also
supports the final Pauli soundness theorem. The dimension-divisibility
remark distinguishes this proved conclusion from the source-law assertions
and records the exceptional defined range $m=1$, $q\ge8$. The soundness
provider remark likewise distinguishes the proved completed-POVM
seed-indexed theorem from the prescribed-effect statement and the printed
import. These remarks record a documented deviation;
their documentary status supplies no formalization mark.
The blueprint's proof node
for the classical soundness import itself states the source's route --- the
identification with the two-prover tensor-code test and the application of
Theorem~4.7 --- and enumerates the two obligations of
Section~\ref{sec:import-obligations} as open, citing this note for the
detailed comparison and the case arithmetic.  The separate support node for
general simultaneity records the established directly indexed theorem and its
proof through the low individual degree theorem. The seed-indexed theorem
for completed point POVMs is proved by compression and global recovery for
arbitrary $k$. The separate node \path{lem:ld-soundness-formalized}
states these completed measurements explicitly and carries both statement
and proof marks. The source-labelled \path{lem:ld-soundness} retains its
prescribed-effect statement and its link to the completed result, but has
neither mark: no linked declaration states the prescribed-effect implication
discussed in Section~2. Its reasoned exemption is recorded in
\doclink{docs/completion/qpbt-leanok-exemptions.md}, under the documentary
authorization for \ghissue{710}. This exemption does not certify that
implication or either printed import obligation.
An earlier revision had
instead added the hypothesis $2m+2 \mid q$ to the three lemmas; that was
statement drift in the sense of the proof-gap protocol, since the hypothesis
excludes all admissible $m \ge 2$ and does not appear in the source.

\section{Conclusion}

The verdict is a \emph{documented deviation}. Direct index sampling is a
justified auxiliary construction: its soundness is proved at every positive
dimension, and it supplies the proved global polynomial-pair conclusion
without $2m+2\mid q$. Seed-indexed soundness for completed point POVMs is
also proved, for every positive simultaneity parameter, on its unchanged
domain $m'\mid q$. Positivity relates this result to the prescribed-effect
conclusion, but the linked Lean theorem does not state that implication.
These proofs use the low individual degree theorem, not the printed
tensor-code import.

The source's balanced seed construction at $2m+2$ remains unavailable when
$m\ge2$ or $(q,m)=(2,1)$. The tensor-code game correspondence and its
parameter substitution remain unverified, and the direct subline law is not
identified with the source law even in the exceptional defined range.
Documentary closure of this register row under \ghissue{710} certifies this
comparison only. It proves no unresolved source-law assertion, refutes no
unknown conclusion, and does not close the separate scalar or error-form
comparisons recorded in Section~\ref{sec:auxiliary-subline-scalars}.

\bibliographystyle{alpha}
\bibliography{references}

\end{document}
